% Riwayat singkat penulis, wajib menurut Permendikbudristek 22/2022 Lampiran III
% Tabel 3. Data Alfa Yohannis dari PDDikti, homepage, dan Google Scholar. Data
% Ariya Panna dari Google Scholar.
\chapter{Riwayat Penulis}

\section*{Alfa Yohannis}

Alfa Yohannis adalah dosen Universitas Pradita di Tangerang. Ia menyelesaikan gelar Doctor of Philosophy bidang ilmu komputer di University of York, Magister Teknik bidang Informatika di Institut Teknologi Bandung, dan Sarjana Teknik bidang Teknik Elektro di Universitas Diponegoro. Sebelumnya ia menjadi \textit{research fellow} di University of York, KTP \textit{associate} di Leonardo Ltd, dan dosen di Kalbis Institute.

Bidang keahliannya meliputi rekayasa perangkat lunak berbasis model, sistem informasi, gamifikasi, sistem terdesentralisasi, dan pendidikan ilmu komputer.

Karya terbitannya antara lain \textit{Turning Models Inside Out} (2017), \textit{Towards Model-Based Bias Mitigation in Machine Learning} (2022), dan \textit{LLM-Assisted TOGAF Preliminary Phase for Student Digital Twins} (2025). Penelitian terakhirnya mengukur jejak karbon dan kinerja penyimpanan data citra deret waktu pada PostgreSQL, MongoDB, dan MinIO.

\section*{Ariya Panna}

Ariya Panna adalah mahasiswa Universitas Pradita. Bidang minatnya meliputi \textit{Internet of Things}, rekayasa \textit{back-end}, dan analisis data. Alamat surelnya \path{ariya.panna@student.pradita.ac.id}.

Karya terbitannya antara lain \textit{Sound-Based Smart Toddler Monitoring System: AIoT Development with YAMNet on Raspberry Pi} (2025), \textit{LLM-Assisted TOGAF Preliminary Phase for Student Digital Twins} (2025), dan \textit{Analisis Perbandingan Kinerja gRPC Go dan Node.js pada Layanan Microtransactions Berbasis Arsitektur Microservices} (2026). Penelitian terakhirnya membandingkan kinerja gRPC pada Go dan Node.js untuk layanan \textit{microtransactions}.
